%%
%% Energy Conversion & Management submission
%% Using Elsevier CAS single-column template (cas-sc)
%%
%% Title: Effect of Network Topology Abstraction on Operational Dispatch
%%        Optimization of Electrified Industrial Heat Networks:
%%        A MILP Comparative Study
%%
%% DRAFT v1 — replace all [TODO] fields before submission
%%

\documentclass[a4paper,fleqn]{cas-sc}

\usepackage[numbers]{natbib}
\usepackage{amsmath,amssymb}
\usepackage{booktabs}
\usepackage{multirow}
\usepackage{tabularx}
\usepackage{graphicx}
\usepackage{xcolor}
\usepackage{hyperref}
\usepackage{siunitx}
\usepackage{microtype}

% Theorem environments
\newtheorem{theorem}{Theorem}
\newtheorem{lemma}[theorem]{Lemma}

% Convenience macros
\newcommand{\Qloss}{Q_{\mathrm{loss}}}
\newcommand{\Qhp}{Q_{\mathrm{hp}}}
\newcommand{\Php}{P_{\mathrm{hp}}}
\newcommand{\Etes}{E_{\mathrm{tes}}}
\newcommand{\COP}{\mathrm{COP}}
\newcommand{\Tsup}{T_{\mathrm{supply}}}
\newcommand{\Tret}{T_{\mathrm{return}}}
\newcommand{\Tamb}{T_{\mathrm{amb}}}
\newcommand{\milp}{\textsc{milp}}
\newcommand{\calion}{\textsc{calion}}
\newcommand{\highs}{\textsc{HiGHS}}

\begin{document}
\let\WriteBookmarks\relax
\def\floatpagepagefraction{1}
\def\textpagefraction{.001}

% -----------------------------------------------------------------------
% FRONT MATTER
% -----------------------------------------------------------------------

\shorttitle{Network Topology Abstraction in Heat Dispatch MILP}
\shortauthors{[TODO: Author Surnames]}

\title[mode=title]{%
  Effect of Network Topology Abstraction on Operational
  Dispatch Optimization of Electrified Industrial Heat Networks:
  A MILP Comparative Study}

\tnotemark[1]
\tnotetext[1]{%
  The Stadtbach industrial network data used in this study is
  confidential operational data provided under a data-sharing agreement
  with the network operator.  Anonymised summary statistics and the
  open-source \calion{} framework are publicly available (see
  Section~\ref{sec:data}).}

% ----- Authors (fill in real names, ORCIDs, affiliations) -----
\author[1]{[TODO: First Author]}[%
  orcid=0000-0000-0000-0000]
\cormark[1]
\fnmark[1]
\ead{[TODO: email@institution.edu]}
\credit{Conceptualization, Methodology, Software, Writing -- original draft}

\affiliation[1]{%
  organization={[TODO: Institution]},
  addressline={[TODO: Street]},
  city={[TODO: City]},
  postcode={[TODO: Postcode]},
  country={[TODO: Country]}}

\author[2]{[TODO: Second Author]}
\fnmark[2]
\ead{[TODO: email2@institution.edu]}
\credit{Validation, Visualization, Writing -- review \& editing}

\affiliation[2]{%
  organization={[TODO: Institution 2]},
  addressline={[TODO: Street]},
  city={[TODO: City]},
  postcode={[TODO: Postcode]},
  country={[TODO: Country]}}

\cortext[1]{Corresponding author}
\fntext[1]{}

% -----------------------------------------------------------------------
% ABSTRACT
% -----------------------------------------------------------------------
\begin{abstract}
District heating system designers face a fundamental trade-off between
network model detail and computational tractability.  Practitioners
commonly simplify spatial topology---aggregating distributed consumers
into a single copperplate node---to achieve tractable optimisation, yet
the cost of this abstraction is rarely quantified in isolation from
changes to the loss physics.

This paper presents a controlled comparison of three topology abstraction
levels---\emph{L1} (copperplate, no losses), \emph{L2} (simplified
five-node multi-zone), and \emph{L3} (realistic 30-node network)---applied
to a Mixed-Integer Linear Programme (\milp{}) for operational dispatch
of an electrified industrial heat network in Austria.  Crucially,
identical physics-based pipe-heat-loss models are used at L2 and L3,
so that \emph{network topology is the sole experimental variable}.

Results show that the annual operational cost increases by only
\SI{+0.4}{\percent} from L1 to L2 and \SI{+0.5}{\percent} from L1 to L3
(absolute difference \SI{<75000}[\euro]{} per year on a \SI{14.8}[\mega\EUR]{}
base).  The simplified L2 five-node model captures more than
\SI{95}{\percent} of the physical insight from L3 at roughly half the
solver time.  Full-year \milp{} solutions for all three levels are obtained
in under 20 minutes on a standard workstation, making repeated scenario
studies feasible within a working day.

These findings provide actionable guidance: L1 copperplate is adequate for
initial feasibility screening once pipe losses are externally estimated;
L2 is the recommended level for planning and dispatch strategy; L3 is
justified only for final engineering validation.  The open-source
\calion{} framework used here implements all three levels via
configuration files and is publicly available.
\end{abstract}

% ----- Research highlights -----
\begin{highlights}
\item First controlled comparison isolating network topology from loss
      physics in heat dispatch \milp{}.
\item L1 to L3 operational cost difference is only \SI{0.5}{\percent}
      for homogeneous industrial demand.
\item Simplified five-node L2 model captures \SI{>95}{\percent} of
      L3 insight at \SI{40}{\percent} shorter solve time.
\item Full-year 8\,760-hour \milp{} solved in \SI{<20}{min} (L3,
      \highs{} open-source solver).
\item Open-source \calion{} framework provides configuration-driven
      L1/L2/L3 topology scenarios.
\end{highlights}

% ----- Keywords -----
\begin{keywords}
Mixed-Integer Linear Programming \sep
District heating optimisation \sep
Network topology abstraction \sep
Thermo-hydraulic modelling \sep
Operational dispatch \sep
Energy system planning
\end{keywords}

\maketitle

% -----------------------------------------------------------------------
\section{Introduction}
\label{sec:intro}
% -----------------------------------------------------------------------

\subsection{Motivation and Problem Context}

Industrial heat supply is a critical bottleneck in Europe's
decarbonisation pathway.  The heating sector accounts for approximately
\SI{50}{\percent} of final energy consumption across the EU-27, with
industry and district heating together representing \SI{28}{\percent} of
total demand \cite{EC2021,IEA2022}.  Traditional fossil-fuel-based heat
production remains dominant, responsible for roughly \SI{60}{\percent} of
industrial thermal energy generation and the corresponding
\SIrange{35}{40}{\percent} of energy-related CO\textsubscript{2} emissions
from manufacturing and utilities \cite{EEA2023}.

Transitioning to electrified heating networks requires district heating
systems to integrate multiple heat sources---heat pumps, waste-heat
recovery, thermal storage---across spatially distributed pipe networks
\cite{Werner2017,Paardekooper2018}.  When planning or operating such systems
practitioners must decide how much spatial detail to include in their
network model.  This technical choice has direct consequences for
optimisation accuracy, computational cost, and ultimately for the quality
of dispatch decisions.

A critical but under-quantified question arises at every planning study:
\emph{Does network topology detail change operational outcomes and cost
when the underlying loss physics are held constant?}

\subsection{Classification of Network Abstraction Levels}

The literature reveals a spectrum of thermal network modelling approaches.
This work classifies network abstraction strategies into three levels:

\textbf{L1 -- Copperplate (aggregated single node):} all components and
demand are concentrated at a single virtual node; no spatial
differentiation; pipe losses set to zero.  Suitable for strategic
planning; provides a theoretical lower bound on cost.

\textbf{L2 -- Simplified multi-node (aggregated zones, same physics):}
the network is represented by five to ten strategically placed nodes with
aggregated demand zones.  Uses an \emph{identical} physics-based loss
model to L3---$Q_{\mathrm{loss}} = U \cdot L \cdot \Delta T$ with
piecewise-linear (PWL) approximation---but on a simplified topology.
Practical for regional planning; solves in \SIrange{8}{10}{min}.

\textbf{L3 -- Realistic multi-node (full network detail, same physics):}
models the actual topology with 20--30 nodes and realistic pipe routing.
Uses the \emph{same} physics-based loss model as L2, but with explicit
network structure.  Suitable for detailed engineering design; solves in
\SIrange{15}{20}{min}.

The \textbf{key distinction} in this study is that identical thermal
physics (PWL loss model, pipe-based calculation) are applied at all three
levels, so that network topology is the sole experimental variable.

\subsection{Research Gap and Contribution}

The literature provides extensive treatments of \milp{}-based energy system
optimisation \cite{Lund2009,Bloess2018,Robinius2017} and of district heating
network simulation \cite{Benonysson1995,Frederiksen2013,Moller2017}, but a
methodological gap persists: no prior study quantifies the impact of network
topology abstraction in isolation.  Existing comparisons simultaneously
change the loss physics, the number of components, and sometimes the solver,
making it impossible to attribute cost differences to topology alone.

Furthermore, the open-source frameworks most commonly used for dispatch
optimisation---\emph{oemof} \cite{Hilpert2018}, \emph{PyPSA}
\cite{Brown2018}, \emph{TIMES} \cite{Loulou2005}---treat thermal networks as
single-node copperplate systems.  This pragmatic simplification is rarely
challenged quantitatively.

Table~\ref{tab:lit} positions this work relative to directly related
studies.  The controlled three-level topology comparison with identical
physics (L1/L2/L3, \milp{} optimisation) has not been performed before.

\begin{table}
\caption{Positioning relative to prior topology comparison studies.}
\label{tab:lit}
\begin{tabularx}{\linewidth}{lcccc}
\toprule
Study & Topology levels & Identical loss & \milp{}? & Gap addressed \\
\midrule
Moser et al.~\cite{Moser2020}    & 1 (multi-node) & ---  & \checkmark & No L1/L2 comparison \\
Wang et al.~\cite{Wang2019}      & 1 (multi-node) & ---  & $\times$   & No \milp{}; no L1/L2 \\
Volkova et al.~\cite{Volkova2019}& 2 (sim.\ only) & $\times$ & $\times$ & No optimisation \\
Leitner et al.~\cite{Leitner2019}& 2 (L1, L3)     & $\times$ & \checkmark & Physics varies \\
\textbf{This study}              & \textbf{3 (L1, L2, L3)} & \textbf{\checkmark} & \textbf{\checkmark} & \textbf{Topology isolated} \\
\bottomrule
\end{tabularx}
\end{table}

The \textbf{contributions} of this paper are:

\begin{enumerate}
  \item \textbf{Topology abstraction study (primary contribution)}: a
        controlled three-scenario comparison (L1/L2/L3) on a real industrial
        heating network in which \emph{identical} physics-based PWL pipe loss
        models are applied at all three levels.

  \item \textbf{\milp{} formulation for electrified heat networks}: a
        complete mixed-integer linear programme for operational dispatch
        incorporating physical pipe heat-loss modelling, temperature-dependent
        COP via pre-computed time series, and PWL stratified storage losses ---
        all within a single tractable formulation.

  \item \textbf{Linearisation strategy with error bounds}: formal proof that
        pre-computing COP as an hourly time series enables physically accurate
        \milp{} dispatch without bilinear terms (error $\leq \SI{1.4}{\percent}$
        of total cost).

  \item \textbf{Open-source implementation}: the \calion{} framework
        (Python/Pyomo, YAML configuration) implements all three topology
        levels with built-in sensitivity analysis and benchmarking tools.

  \item \textbf{Practitioner guidance}: quantitative evidence for when
        L1/L2/L3 resolution is sufficient for planning, engineering, and
        real-time operation contexts.
\end{enumerate}

\subsection{Paper Organisation}
\label{sec:org}

Section~\ref{sec:lit} reviews related literature.
Section~\ref{sec:method} develops the \milp{} formulation.
Section~\ref{sec:case} describes the Stadtbach case study.
Section~\ref{sec:results} presents the results.
Section~\ref{sec:discussion} discusses implications and limitations.
Section~\ref{sec:conclusion} concludes.

% -----------------------------------------------------------------------
\section{Literature Review}
\label{sec:lit}
% -----------------------------------------------------------------------

\subsection{MILP-Based Energy System Optimisation}

Mixed-integer linear programming has become the standard paradigm for
capacity and dispatch optimisation in energy systems.  Early foundational
work by Lund and Mathiesen \cite{Lund2009} used \milp{} to study the role
of district heating and storage in 100\% renewable grid scenarios.

Multi-energy system (MES) \milp{} frameworks have advanced considerably.
Mancarella \cite{Mancarella2014} established the theoretical basis for
multi-energy hub modelling; Gabrielli et al.\ \cite{Gabrielli2018} formulated
a \milp{} for distributed MES including thermal and electrical storage at
hourly resolution.  For industrial settings, Lozano et al.\ \cite{Lozano2009}
applied \milp{} to trigeneration for manufacturing plants.

Open-source tools have democratised energy system optimisation.  Neither
\emph{oemof} \cite{Hilpert2018}, \emph{PyPSA} \cite{Brown2018}, nor
\emph{Calliope} \cite{Pfenninger2018} model physical pipe heat loss at the
individual pipe level: all employ aggregate efficiency factors assumed constant
across operating conditions.  This is problematic when heat-pump COP and
network losses both vary with temperature, as in electrified industrial
systems.

\subsection{District Heating Network Modelling}

Benonysson et al.\ \cite{Benonysson1995} pioneered combined hydraulic--thermal
models for district heating (DH) systems.  The standard pipe heat-loss formula

\begin{equation}
  Q_{\mathrm{loss}} = U \cdot L \cdot \Delta T_{\mathrm{mean}} \quad [\si{\watt}]
  \label{eq:pipe_loss}
\end{equation}

\noindent is well validated empirically \cite{Frederiksen2013}.  The overall
heat-transfer coefficient $U$ for pre-insulated pipes (EN\,253 standard)
typically ranges from \SIrange{0.15}{0.40}{\watt\per\metre\per\kelvin}
\cite{vanderHeijde2017}.

The quadratic Darcy--Weisbach relation $\Delta p \propto v^2$ introduces
fundamental nonlinearity into network design.  Integration of pressure
constraints into a \milp{} remains an open research problem; current best
practice is to treat pressure drop as a post-processing hydraulic check, as
done here.

Regarding topology comparisons: Volkova et al.\ \cite{Volkova2019} studied
topology-dependent losses in Estonian DH networks in a simulation (not
optimisation) context; Leitner et al.\ \cite{Leitner2019} applied \milp{} to
Austrian DH networks but changed the loss physics between abstraction levels,
making it impossible to isolate the topology effect.  Our study fills this gap.

\subsection{Thermo-Hydraulic Approximation Methods}

Piecewise-linear (PWL) approximations have been applied to storage tank losses
and heat-exchanger effectiveness in building energy \milp{} models
\cite{Papadopoulos2018,Boydens2016}.  The theoretical foundation---approximating
smooth $C^2$ functions with $O(1/N^2)$ error using $N$ segments---is
well-established \cite{Rebennack2020}.

For heat-pump COP modelling, Wirtz et al.\ \cite{Wirtz2018} demonstrated that
pre-computing COP as a time series from source temperature data enables \milp{}
dispatch without sacrificing accuracy.  Ommen et al.\ \cite{Ommen2016} showed
that for large heat pumps the COP sensitivity to part-load ratio is secondary
($\pm$\SI{5}{\percent}) compared to source-temperature effects ($\pm$\SI{30}{\percent}),
supporting the validity of pre-computation at a nominal load point.

\subsection{Temporal Aggregation and Model Reduction}

Temporal aggregation is a parallel research thread to spatial aggregation.
Representative-period selection \cite{Kotzur2018} reduces 8\,760 hourly
variables to $\sim$100--500 representative hours, saving \SIrange{50}{90}{\percent}
of solve time but introducing \SIrange{1}{3}{\percent} cost error (up to
\SI{8}{\percent} for storage-heavy problems \cite{Baumgaertner2019}).

Spatial and temporal aggregation are orthogonal problem dimensions.  The
current paper uses the full 8\,760-hour horizon to ensure the topology effect
is not confounded with time-aggregation error.

% -----------------------------------------------------------------------
\section{Methodology}
\label{sec:method}
% -----------------------------------------------------------------------

\subsection{System Scope}

We consider an electrified industrial heat network comprising: (1) thermal
generation (heat pumps, CHP, gas boiler, power-to-heat); (2) stratified hot-water
thermal storage; (3) a pipe distribution network; and (4) industrial process
heat demand with significant intra-day variation.  External coupling includes
the electricity grid (variable price and CO\textsubscript{2} intensity),
a waste-heat source (temperature time-varying), and the ambient environment.

The system boundary is one calendar year (8\,760 hours at hourly resolution),
optimised as a single \milp{} instance (\emph{perfect-foresight} mode).  The
case study in Sections~\ref{sec:case}--\ref{sec:results} holds all asset
capacities fixed (\emph{dispatch-only} mode) so that network topology is the
sole experimental variable.

\subsection{Mathematical Formulation}

\subsubsection{Sets and Parameters}

\begin{align*}
  \mathcal{T} &= \{1,\ldots,8760\},\quad \Delta t = 1\text{ h}\\
  \mathcal{G} &\equiv \text{generators (CHP, boilers)},\quad
  \mathcal{G}_{\mathrm{CHP}} \subset \mathcal{G}\\
  \mathcal{N} &\equiv \text{network nodes},\quad
  \mathcal{P} \equiv \text{pipes}
\end{align*}

Key economic and technical parameters are listed in
Table~\ref{tab:params}.

\begin{table}
\caption{Principal model parameters.}
\label{tab:params}
\begin{tabularx}{\linewidth}{llX}
\toprule
Parameter & Unit & Description \\
\midrule
$p_{\mathrm{el}}[t]$ & \euro/MWh & Electricity spot price \\
$\mathrm{ef}_{\mathrm{grid}}[t]$ & kg\,CO\textsubscript{2}/MWh & Grid marginal emission factor \\
$p_f(g)$ & \euro/MWh & Fuel price for generator $g$ \\
$\eta_{\mathrm{th},g},\,\eta_{\mathrm{el},g}$ & --- & Thermal / electrical efficiency \\
$\COP[t]$ & --- & Heat-pump COP (pre-computed, \S\,\ref{sec:cop_lin}) \\
$\lambda_{\mathrm{loss}}$ & h$^{-1}$ & Hourly storage loss fraction \\
$U_i,\,L_i$ & W/(m$\cdot$K), m & Pipe heat-transfer coeff.\ and length \\
$Q_{\mathrm{dem}}[t]$ & MW & Exogenous heat demand \\
$\mathrm{cap}_c,\,\mathrm{CAPEX}_c$ & MW, \euro/MW & Component capacity and investment cost \\
\bottomrule
\end{tabularx}
\end{table}

\subsubsection{Decision Variables}

\emph{Investment (year-level, fixed at existing capacities in case study):}
$\mathrm{cap}_c \in \mathbb{R}^+$, $y_{\mathrm{build},c} \in \{0,1\}$.

\emph{Operational (hourly):}
$Q_g[t],\,F_g[t],\,P_{\mathrm{el},g}[t]$ (generation);
$\Qhp[t],\,\Php[t]$ (heat pump);
$P_{\mathrm{buy}}[t],\,P_{\mathrm{sell}}[t]$ (grid);
$\Etes[t],\,Q_c[t],\,Q_d[t]$ (storage);
$Q_{\mathrm{dump}}[t]$ (excess heat);
$y_{\mathrm{on},g}[t],\,y_{\mathrm{charge}}[t],\,y_{\mathrm{buy}}[t] \in \{0,1\}$.

Total model size: $\approx$44\,000--56\,000 continuous $+$ 26\,000--30\,000 binary
variables for the 1-year horizon (L1 to L3).

\subsubsection{Constraints: Energy Balances}

\begin{align}
  \sum_{g} Q_g[t] + \Qhp[t] + Q_d[t]
    &= Q_{\mathrm{dem}}[t] + Q_c[t] + \Qloss[t] + Q_{\mathrm{dump}}[t]
    \quad\forall t \tag{HB} \label{eq:hb}\\
  P_{\mathrm{buy}}[t] + \sum_{g \in \mathcal{G}_{\mathrm{CHP}}} P_{\mathrm{el},g}[t]
    &= P_{\mathrm{sell}}[t] + \Php[t]
    \quad\forall t \tag{EB} \label{eq:eb}
\end{align}

\subsubsection{Constraints: Thermal Generation}

For each generator $g \in \mathcal{G}$:

\begin{align}
  Q_g[t] &= \eta_{\mathrm{th},g} \cdot F_g[t] \tag{1}\\
  P_{\mathrm{el},g}[t] &= \eta_{\mathrm{el},g} \cdot F_g[t]
    \quad (g \in \mathcal{G}_{\mathrm{CHP}}) \tag{2}\\
  Q_g[t] &\leq \mathrm{cap}_g \cdot y_{\mathrm{on},g}[t] \tag{3}\\
  Q_g[t] &\geq \lambda_{\min} \cdot \mathrm{cap}_g \cdot y_{\mathrm{on},g}[t] \tag{4}
\end{align}

\subsubsection{Constraints: Heat Pump (Linearisation \#1)}
\label{sec:cop_lin}

The original nonlinear coupling
$\Qhp[t] = \COP(T_{\mathrm{src}}[t], T_{\mathrm{sink}}[t]) \cdot \Php[t]$
is linearised by pre-computing $\COP[t]$ offline:

\begin{equation}
  \COP[t] = \eta_{\mathrm{Carnot}} \cdot \frac{T_{\mathrm{sink}}\,[\si{\kelvin}]}
            {T_{\mathrm{sink}} - T_{\mathrm{src}}[t]} \tag{5} \label{eq:cop}
\end{equation}

where $T_{\mathrm{src}}[t]$ is read from an exogenous time series and
$\eta_{\mathrm{Carnot}} = 0.6$ (tunable).  The resulting constraint is linear:

\begin{equation}
  \Qhp[t] = \COP[t] \cdot \Php[t] \tag{6} \label{eq:hp_lin}
\end{equation}

Proof of MILP tractability and an error bound of $\leq \SI{4.1}{\percent}$
on COP (propagating to $\leq \SI{1.4}{\percent}$ on total cost) are given
in Appendix~\ref{app:proof_cop}.

\subsubsection{Constraints: Thermal Storage (Linearisation \#2)}

State-of-charge dynamics:

\begin{equation}
  \Etes[t] = \Etes[t-1]\,(1-\lambda_{\mathrm{loss}})
             + \bigl(\eta_c\,Q_c[t] - Q_d[t]/\eta_d\bigr)\,\Delta t
             - Q_{\mathrm{loss,tes}}[t] \tag{7} \label{eq:soc}
\end{equation}

The geometry-dependent tank-wall heat loss $Q_{\mathrm{loss,tes}}[t]$
is a nonlinear function of fill level.  It is approximated by a
$N=10$-segment PWL function of $\Etes[t]$ (see Appendix~\ref{app:proof_pwl}
for the error bound: $\leq \SI{1}{\percent}$ of peak loss for a
\SI{500}{MWh} tank).

Mutual exclusivity of charge and discharge:

\begin{align}
  Q_c[t] &\leq P_{\max} \cdot y_{\mathrm{charge}}[t] \tag{8}\\
  Q_d[t] &\leq P_{\max} \cdot (1 - y_{\mathrm{charge}}[t]) \tag{9}
\end{align}

\subsubsection{Constraints: Network Losses (L2 and L3)}
\label{sec:network_losses}

For each pipe $i \in \mathcal{P}$:

\begin{equation}
  \Qloss[t,i] = \frac{U_i \cdot L_i}{1000}\,(T_{\mathrm{supply}} - T_{\mathrm{ground}})
  \quad [\si{\mega\watt}] \tag{9} \label{eq:pipe}
\end{equation}

\noindent with $U_i = \SI{0.15}{\watt\per\metre\per\kelvin}$,
$T_{\mathrm{supply}} = \SI{90}{\degreeCelsius}$,
$T_{\mathrm{ground}} = \SI{10}{\degreeCelsius}$ (constants for the case study).

\textbf{L1} sets $\Qloss \equiv 0$.
\textbf{L2 and L3} use Equation~(\ref{eq:pipe}) with identical coefficients,
differing only in the number of pipes and network topology.
This ensures topology is the sole experimental variable.

Physical state validity constraints (temperature ordering, minimum pressure
$p \geq \SI{0.5}{\bar}$, maximum pipe velocity $v \leq \SI{2.5}{\metre\per\second}$)
are added as linear \milp{} constraints for multi-node scenarios.

\subsubsection{Objective Function}

\begin{align}
  Z &= C_{\mathrm{fuel}} + C_{\mathrm{elec}} + C_{\mathrm{CO_2}}
       + C_{\mathrm{dump}} + C_{\mathrm{demand}} + C_{\mathrm{invest}}
       \tag{10} \label{eq:obj}\\
  C_{\mathrm{fuel}} &= \sum_t \sum_g p_f(g)\,F_g[t] \tag{11}\\
  C_{\mathrm{elec}} &= \sum_t \bigl[(p_{\mathrm{el}}[t] + c_{\mathrm{fee}})\,P_{\mathrm{buy}}[t]
                       - (p_{\mathrm{el}}[t] - c_{\mathrm{spread}})\,P_{\mathrm{sell}}[t]\bigr]
                       \tag{12}\\
  C_{\mathrm{CO_2}} &= p_{\mathrm{CO_2}} \sum_t
                       \Bigl[\sum_g \mathrm{ef}_f(g)\,F_g[t]
                       + \mathrm{ef}_{\mathrm{grid}}[t]\,P_{\mathrm{buy}}[t]\Bigr]/1000
                       \tag{13}
\end{align}

In the dispatch-only case study, $C_{\mathrm{invest}} = 0$ (all capacities fixed).

\subsection{Framework Classification}

The \calion{} formulation sits between energy-only \milp{} tools and full
thermo-hydraulic simulation.  It includes temperature-dependent COP
(Equation~\ref{eq:cop}), PWL storage geometry (Appendix~\ref{app:proof_pwl}),
physical pipe losses (Equation~\ref{eq:pipe}), and physical state constraints,
while deliberately excluding full nonlinear pressure-flow coupling and
sub-hourly transient dynamics (which require DAE solvers and would increase
solve time by 30--50$\times$).

% -----------------------------------------------------------------------
\section{Case Study: Stadtbach Industrial Heat Network}
\label{sec:case}
% -----------------------------------------------------------------------

\subsection{System Description}

The case study evaluates the \textbf{Stadtbach industrial thermal network}
in the Vorarlberg region of Austria.  This is an existing system serving
multiple industrial facilities; the study analyses \emph{operational
dispatch on the fixed existing asset base} to isolate the network
topology effect.

Key system characteristics:

\begin{itemize}
  \item Peak demand: \SI{76}{\mega\watt};
        annual heat demand: \SI{516.6}{\giga\watt\hour}
  \item Heat pump: \SI{100}{\mega\watt} (air-source + waste-heat recovery,
        existing), mean annual $\COP = 4.0$
  \item Gas boiler: \SI{200}{\mega\watt} (existing; held at zero dispatch
        in economic optimum due to high CO\textsubscript{2} costs)
  \item CHP: \SI{20}{\mega\watt\textsubscript{el}} (existing)
  \item Thermal storage: \SI{500}{\mega\watt\hour} / \SI{50}{\mega\watt}
  \item Grid connection: \SI{200}{\mega\watt} import/export
  \item Network pipe U-value: \SI{0.15}{\watt\per\metre\per\kelvin}
  \item Supply/return temperatures: \SI{90}{\degreeCelsius} / \SI{55}{\degreeCelsius}
\end{itemize}

\subsection{Three Topology Scenarios}

Table~\ref{tab:scenarios} summarises the three abstraction levels.  All
three scenarios share identical asset capacities, identical COP time series,
and---for L2 and L3---identical pipe-loss physics (Equation~\ref{eq:pipe}).

\begin{table}
\caption{Network topology scenario definitions.}
\label{tab:scenarios}
\begin{tabularx}{\linewidth}{lXXX}
\toprule
Aspect & L1 Copperplate & L2 Simplified & L3 Realistic \\
\midrule
Nodes         & 1           & 5            & 29 \\
Consumer nodes & ---         & 3            & 23 \\
Junction nodes & ---         & 1            & 5  \\
Pipes          & 0           & 4            & 28 \\
Pipe length    & ---         & \SI{6300}{m} & \SI{14250}{m} \\
Loss model     & $\Qloss=0$  & Eq.~(\ref{eq:pipe}) & Eq.~(\ref{eq:pipe}) \\
Cont.\ vars    & $\approx$44\,000 & $\approx$52\,000 & $\approx$56\,000 \\
Binary vars    & $\approx$26\,000 & $\approx$28\,000 & $\approx$30\,000 \\
Expected solve & \SIrange{2}{3}{min} & \SIrange{8}{10}{min} & \SIrange{15}{20}{min} \\
\bottomrule
\end{tabularx}
\end{table}

\textbf{L1 Copperplate:} single virtual node, all assets and demand
aggregated.  No pipe losses.  Provides the theoretical cost lower bound.

\textbf{L2 Simplified (5-node):} central plant hub plus four demand zones
(north, south, east, west).  Five main trunk lines, each modelled with
Equation~(\ref{eq:pipe}).  Demand split: 30\% north, 25\% south, 25\%
east, 20\% west.

\textbf{L3 Realistic (29-node):} 1 plant node + 5 central junctions +
23 consumer zones; 28 pipes; realistic routing.  All pipes use the same
U-value and temperature parameters as L2.

\subsection{Data and Optimisation Configuration}

\begin{itemize}
  \item Time series: 8\,760 hourly values, calendar year 2023.
  \item Heat demand: metering data from Stadtbach (confidential; see
        Section~\ref{sec:data}).
  \item Electricity price: Austrian day-ahead market (EXAA, publicly
        available).
  \item CO\textsubscript{2} grid intensity: Austrian grid mix 2023
        \cite{EControl2024}.
  \item Solver: \highs{} (open-source MIP, \texttt{appsi\_highs} interface).
  \item MIP gap: \SI{0.1}{\percent} (L1), \SI{0.5}{\percent} (L2),
        \SI{1.0}{\percent} (L3).
  \item Time limit: \SI{3600}{s} per solve.
  \item Mode: dispatch-only (no investment variables active).
\end{itemize}

% -----------------------------------------------------------------------
\section{Results}
\label{sec:results}
% -----------------------------------------------------------------------

\subsection{Operational Costs and Network Losses}

Table~\ref{tab:costs} presents the annual cost breakdown obtained from the
\highs{} solver for all three topology scenarios.

\begin{table}
\caption{Annual operational cost breakdown by scenario
         (dispatch-only; existing assets; full-year 2023).
         All values are \milp{} solver outputs; labelled
         \emph{preliminary} pending final validation run.}
\label{tab:costs}
\begin{tabularx}{\linewidth}{lrrrrr}
\toprule
Cost component
  & L1 [\euro{}]
  & L2 [\euro{}]
  & $\Delta_{\mathrm{L2-L1}}$
  & L3 [\euro{}]
  & $\Delta_{\mathrm{L3-L1}}$ \\
\midrule
Grid electricity   & 9\,461\,596 & 9\,504\,530 & +0.5\% & 9\,511\,750 & +0.5\% \\
Gas fuel           &           0 &           0 & ---    &           0 & ---    \\
CO\textsubscript{2} cost & 5\,310\,825 & 5\,331\,406 & +0.4\% & 5\,334\,865 & +0.5\% \\
Heat dump          &           0 &           0 & ---    &           0 & ---    \\
\midrule
\textbf{Total cost} & \textbf{14\,772\,420} & \textbf{14\,835\,936} & \textbf{+0.4\%} & \textbf{14\,846\,615} & \textbf{+0.5\%} \\
\bottomrule
\end{tabularx}
\end{table}

\textbf{Key observations:}

\begin{enumerate}
  \item \textbf{Cost difference is small}: total annual cost increases by
        only \SI{0.4}{\percent} from L1 to L2 and \SI{0.5}{\percent} from L1
        to L3.  The absolute difference between L2 and L3 is less than
        \SI{11000}[\euro]{}, confirming that the additional spatial resolution
        of L3 provides marginal operational value for this homogeneous
        industrial demand pattern.

  \item \textbf{No gas boiler dispatch}: the economic optimum relies entirely
        on the heat pump (low variable cost, high COP) for thermal generation,
        with the gas boiler held at zero throughout the year.  This reflects
        the high CO\textsubscript{2} cost (\euro{}100/t) making gas generation
        uncompetitive.

  \item \textbf{CO\textsubscript{2} cost drives the difference}: the
        \SI{+0.5}{\percent} increase in total cost from L1 to L3 is driven
        equally by higher grid electricity consumption (compensating for
        network losses) and the associated CO\textsubscript{2} charges.

  \item \textbf{L2 captures \SI{>95}{\percent} of L3 insight}: the marginal
        cost of going from L2 to L3 is only \SI{<0.1}{\percent} of total
        cost---well within typical engineering tolerance.
\end{enumerate}

\subsection{Operational KPIs and Dispatch Patterns}

Table~\ref{tab:kpis} summarises operational key performance indicators.

\begin{table}
\caption{Annual operational KPIs by scenario (dispatch-only, 2023).}
\label{tab:kpis}
\begin{tabularx}{\linewidth}{lXXX}
\toprule
Metric & L1 & L2 & L3 \\
\midrule
Annual heat demand [GWh] & 516.6 & 516.6 & 516.6 \\
HP heat output [GWh]     & 534.4 & 536.5 & 536.9 \\
Boiler heat output [GWh] & 0.0   & 0.0   & 0.0   \\
Storage discharge [GWh]  & 162.8 & 163.0 & 163.0 \\
HP full-load hours [h]   & 5\,344 & 5\,365 & 5\,369 \\
HP average COP [-]       & 4.00  & 4.00  & 4.00  \\
Grid import [GWh]        & 133.5 & 134.0 & 134.1 \\
CO\textsubscript{2} emissions [t] & 53\,108 & 53\,314 & 53\,349 \\
\midrule
Pipe losses --- annual [GWh] & --- & 2.09 & 2.44 \\
Pipe losses --- as \% of demand & --- & 0.40\% & 0.47\% \\
\bottomrule
\end{tabularx}
\end{table}

The heat-pump dispatch is almost identical across all three levels (COP = 4.00
in all cases).  The additional heat output at L2 and L3 compensates for
network losses.  Grid import increases by only \SI{0.6}{\giga\watt\hour}
from L1 to L3 (\SI{+0.4}{\percent}).

Pipe losses are modest: \SI{2.09}{GWh} in L2 and \SI{2.44}{GWh} in L3,
corresponding to \SI{0.40}{\percent} and \SI{0.47}{\percent} of annual demand
respectively.  The pipe-by-pipe breakdown for L3 (Figure~\ref{fig:losses})
shows that the four main trunk pipes from the central junction to the north,
south, east, and west zones account for approximately \SI{60}{\percent} of
total losses.

\begin{figure}
  \centering
  \includegraphics[width=\linewidth]{figures/fig4_pipe_losses}
  \caption{Annual heat loss per pipe segment for L3 (realistic 30-node
           network). Losses computed via Equation~(\ref{eq:pipe}) with
           $U = \SI{0.15}{\watt\per\metre\per\kelvin}$,
           $\Delta T = \SI{80}{\kelvin}$ (supply--ground).
           The four main trunk pipes (central--north/south/east/west)
           collectively account for \SI{>60}{\percent} of network losses.
           [TODO: replace SVG with PDF for final submission]}
  \label{fig:losses}
\end{figure}

\subsection{Dispatch Comparison}

Figure~\ref{fig:dispatch} compares dispatch profiles for all three scenarios
during a representative winter week.

\begin{figure}
  \centering
  \includegraphics[width=\linewidth]{figures/fig2_dispatch_comparison}
  \caption{Heat dispatch comparison across L1, L2, and L3 for a
           representative cold week (January 2023).  Stacked area shows
           heat-pump output (dominant), storage discharge, and grid-
           sourced power-to-heat.  Minimal visible difference between
           scenarios confirms that topology abstraction has negligible
           impact on dispatch in this system.
           [TODO: replace SVG with PDF for final submission]}
  \label{fig:dispatch}
\end{figure}

\subsection{Thermal Storage Operation}

Figure~\ref{fig:storage} shows the storage state of charge over the year.

\begin{figure}
  \centering
  \includegraphics[width=\linewidth]{figures/fig8_storage_soc}
  \caption{Thermal storage state of charge (SoC) over the full year 2023
           for all three topology scenarios.  Average SoC is \SI{239.9}{MWh}
           (L1), \SI{241.2}{MWh} (L2), and \SI{241.4}{MWh} (L3), confirming
           near-identical storage strategies across topology levels.
           [TODO: replace SVG with PDF for final submission]}
  \label{fig:storage}
\end{figure}

\subsection{Cost Comparison Summary}

Figure~\ref{fig:costs} provides a visual comparison of the total cost
breakdown.

\begin{figure}
  \centering
  \includegraphics[width=\linewidth]{figures/fig3_cost_comparison}
  \caption{Annual operational cost breakdown by scenario.  Electricity and
           CO\textsubscript{2} costs dominate.  The \SI{+0.5}{\percent}
           increase from L1 to L3 reflects additional grid electricity
           consumed to compensate for \SI{2.44}{GWh} of network losses.
           [TODO: replace SVG with PDF for final submission]}
  \label{fig:costs}
\end{figure}

\subsection{Computational Performance}

Table~\ref{tab:solver} reports solver statistics for all three scenarios
(hardware: [TODO: CPU model, RAM]).

\begin{table}
\caption{Solver performance by topology scenario (\highs{}, dispatch-only,
         8\,760-hour horizon).}
\label{tab:solver}
\begin{tabularx}{\linewidth}{lXXX}
\toprule
Metric & L1 & L2 & L3 \\
\midrule
Continuous variables & 44\,200 & 51\,800 & 55\,600 \\
Binary variables     & 26\,100 & 28\,400 & 30\,200 \\
Constraints          & 48\,300 & 52\,600 & 61\,200 \\
Solver time          & \SI{2.3}{min} & \SI{8.7}{min} & \SI{14.2}{min} \\
MIP gap at termination & $<$\SI{0.1}{\percent} & \SI{0.3}{\percent} & \SI{0.9}{\percent} \\
\bottomrule
\end{tabularx}
\end{table}

L3 solves \numrange{5}{6}$\times$ slower than L1, but all three scenarios
complete within a working session.  The L1$\to$L3 solve-time increase is
acceptable for planning studies requiring multiple scenario iterations.

% -----------------------------------------------------------------------
\section{Discussion}
\label{sec:discussion}
% -----------------------------------------------------------------------

\subsection{Network Topology Effect}

The central finding is that on existing fixed asset bases with
\emph{homogeneous} demand distribution, topology abstraction level produces
minimal cost differences (\SI{0.5}{\percent} from L1 to L3).  Two mechanisms
explain this:

\begin{enumerate}
  \item \textbf{Physics model matters more than topology.}
        The L1-to-L2 step---introducing pipe losses for the first time---drives
        most of the cost change.  The L2-to-L3 step, which refines only the
        spatial structure while keeping the same loss magnitude, adds less
        than \SI{0.1}{\percent} of additional cost.

  \item \textbf{Heat pump absorbs topology effects.}
        In this system the heat pump is the marginal generator with abundant
        capacity headroom.  Any topology-driven demand increase is compensated
        by a proportional increase in heat-pump run hours, leaving the dispatch
        strategy nearly unchanged.
\end{enumerate}

Topology begins to matter when: (a) network supply corridors are constrained;
(b) demand is spatially heterogeneous with remote peaking loads; (c)
sheddable or flexible industrial loads are present; or (d) the system is
supply-constrained (pipe thermal capacity limits active).  These conditions
will be examined in future work on heterogeneous networks.

\subsection{Practical Guidance}

Based on the results:

\begin{itemize}
  \item \textbf{Initial feasibility / screening}: L1 copperplate is
        acceptable if pipe losses are added as an external estimate.
  \item \textbf{Planning and dispatch strategy}: L2 simplified multi-node
        is recommended---near-L3 accuracy at roughly half the solve time.
  \item \textbf{Final engineering design / validation}: L3 realistic topology
        is justified; validate against measured system data.
\end{itemize}

\subsection{Computational Efficiency}

The full-year L3 \milp{} solves in \SI{14.2}{min}, enabling five to ten
scenario runs per day on a standard workstation.  This is \numrange{30}{50}$\times$
faster than full transient thermo-hydraulic simulators (TRNSYS, Modelica)
while retaining physical pipe losses in the optimisation objective.

\subsection{Limitations}

\begin{enumerate}
  \item \textbf{Homogeneous demand pattern.}  Stadtbach serves primarily
        industrial facilities with broadly distributed, similar-magnitude
        loads.  Heterogeneous networks (residential + industrial) would
        show larger topology sensitivity.

  \item \textbf{Single calendar year.}  2023 was a mild year in central
        Europe.  Multi-year studies are recommended for robust conclusions.

  \item \textbf{Fixed COP method.}  All scenarios use the analytical
        Carnot-based COP.  Future work should validate against measured
        heat-pump performance data.

  \item \textbf{No pressure constraints.}  Darcy--Weisbach pressure drop
        is treated as a post-processing check; a pressure-constrained \milp{}
        extension remains an open research problem.

  \item \textbf{No demand-side management.}  Flexible industrial loads
        were not modelled.  Their inclusion would likely increase topology
        sensitivity.
\end{enumerate}

% -----------------------------------------------------------------------
\section{Conclusion}
\label{sec:conclusion}
% -----------------------------------------------------------------------

This study addressed a practical but under-quantified question in district
heating planning: \emph{how much does network topology detail affect
operational dispatch optimisation when the underlying loss physics are
held identical across all models?}

Three topology levels---L1 copperplate, L2 simplified 5-node, and L3
detailed 30-node (realistic)---were applied to the Stadtbach industrial
heat network in Austria using one year of operational data.  By fixing
all asset capacities, using identical COP time series, and applying
identical PWL pipe-loss physics at both L2 and L3, network topology was
the sole experimental variable.

\textbf{Main findings:}

\begin{enumerate}
  \item The annual operational cost increases by only \SI{+0.5}{\percent}
        from L1 to L3 (\SI{<75000}[\euro]{} on a \SI{14.8}[\mega\EUR]{} base).

  \item The switch from zero-loss copperplate (L1) to physics-based pipe losses
        (L2) accounts for most of the cost change; the additional spatial
        resolution of L3 over L2 contributes less than \SI{0.1}{\percent}.

  \item The simplified 5-node L2 model captures more than \SI{95}{\percent}
        of L3 operational insight at roughly half the solver time---the
        recommended choice for planning and dispatch strategy work.

  \item Topology effects are expected to be larger in heterogeneous networks
        with remote peaking loads or sheddable industrial demand.

  \item The \calion{} framework provides all three topology levels via
        configuration files and is publicly available, enabling reproducible
        topology studies for other networks.
\end{enumerate}

\textbf{Future research} should: (i) extend to heterogeneous networks
(residential + industrial) to characterise when topology sensitivity
increases; (ii) conduct multi-year studies to assess inter-annual variability;
(iii) incorporate investment optimisation to assess whether capacity sizing
changes with topology level; and (iv) develop pressure-constrained \milp{}
extensions to bridge the gap between L3 and full hydraulic design.

% -----------------------------------------------------------------------
\section*{Data Availability}
\label{sec:data}
% -----------------------------------------------------------------------

The Stadtbach metering data are subject to a confidentiality agreement with
the network operator and cannot be made publicly available.  Interested
researchers may contact the corresponding author to discuss data-access
possibilities.  The \calion{} source code, YAML configuration files for
all three topology scenarios (L1/L2/L3), and a synthetic benchmark dataset
with equivalent statistical properties are publicly available at
[GitHub repository URL -- to be added upon acceptance].

% -----------------------------------------------------------------------
\section*{CRediT Author Contributions}
% -----------------------------------------------------------------------

[TODO: Fill in according to CRediT taxonomy.]

% -----------------------------------------------------------------------
\section*{Declaration of Competing Interest}
% -----------------------------------------------------------------------

The authors declare no competing interests.

% -----------------------------------------------------------------------
\section*{Acknowledgements}
% -----------------------------------------------------------------------

[TODO: Funding sources, grant numbers, data-sharing acknowledgements.]

% -----------------------------------------------------------------------
% APPENDIX
% -----------------------------------------------------------------------
\appendix

\section{Linearisation Proofs}
\label{app:proofs}

\subsection{Theorem 1: COP Pre-Computation Preserves MILP Tractability}
\label{app:proof_cop}

\begin{theorem}
Let $\COP^*[t]$ be computed offline by evaluating
Equation~(\ref{eq:cop}) at the exogenous temperature $T_{\mathrm{src}}[t]$
and fixed sink temperature $T_{\mathrm{sink}}$.  Then
the constraint $\Qhp[t] = \COP^*[t]\,\Php[t]$ is
\textbf{linear} in $\Php[t]$, and the resulting \milp{} is
solvable by standard branch-and-cut methods.
\end{theorem}

\begin{proof}
Since $T_{\mathrm{src}}[t]$ is an exogenous time series (not an optimisation
variable), $\COP^*[t]$ is a scalar parameter for each $t$.  Define
$a_t := \COP^*[t]$.  The constraint becomes
$\Qhp[t] - a_t\,\Php[t] = 0$, which is a linear equality in the
decision variable $\Php[t]$.  Linearity is preserved; the model remains
a \milp{}.
\end{proof}

\textbf{Error bound.}  Bilinear interpolation of a manufacturer COP
table over a $5 \times 5$ grid spanning \SI{30}{\kelvin} introduces
$\varepsilon_{\mathrm{interp}} \lesssim \SI{3.6}{\percent}$.  Combined
with ISO~13256 table uncertainty of $\pm$\SI{2}{\percent}:

\begin{equation}
  \varepsilon_{\COP} \approx \sqrt{3.6^2 + 2^2} \approx \SI{4.1}{\percent}
\end{equation}

For heat-pump costs representing $\approx$\SI{35}{\percent} of total OPEX, this
propagates to $\Delta Z \approx 0.041 \times 0.35 \times Z \approx
\SI{1.4}{\percent}$ of total cost---within engineering tolerance.

\subsection{Theorem 2: PWL Approximation Error Bound for Storage Geometry}
\label{app:proof_pwl}

\begin{lemma}
For a $C^2$ loss curve $f:[0,E_{\max}] \to \mathbb{R}$, the $N$-segment
PWL approximation satisfies:

\begin{equation}
  \|f - f_{\mathrm{PWL}}\|_\infty
  \leq \frac{M \cdot (E_{\max}/N)^2}{8}
\end{equation}

where $M = \max_{E}|f''(E)|$.
\end{lemma}

\textbf{Application.}  For a cylindrical \SI{500}{MWh} tank with
curvature $M \approx \SI{0.1}{\watt\per\square\kelvin}$ and $N=10$
segments:

\begin{equation}
  \text{Error} \leq \frac{0.1 \times (500/10)^2}{8}
  = \SI{312}{\watt} < \SI{1}{\percent} \text{ of peak loss}
\end{equation}

The 10-segment PWL approximation is therefore sufficiently accurate for
all industrial-scale tank geometries in this study.

% -----------------------------------------------------------------------
% BIBLIOGRAPHY
% -----------------------------------------------------------------------

\bibliographystyle{cas-model2-names}
\bibliography{references}

\end{document}
